% Options for packages loaded elsewhere
\PassOptionsToPackage{unicode}{hyperref}
\PassOptionsToPackage{hyphens}{url}
\documentclass[
  12pt,
]{article}
\usepackage{xcolor}
\usepackage{amsmath,amssymb}
\setcounter{secnumdepth}{-\maxdimen} % remove section numbering
\usepackage{iftex}
\ifPDFTeX
  \usepackage[T1]{fontenc}
  \usepackage[utf8]{inputenc}
  \usepackage{textcomp} % provide euro and other symbols
\else % if luatex or xetex
  \usepackage{unicode-math} % this also loads fontspec
  \defaultfontfeatures{Scale=MatchLowercase}
  \defaultfontfeatures[\rmfamily]{Ligatures=TeX,Scale=1}
\fi
\usepackage{lmodern}
\ifPDFTeX\else
  % xetex/luatex font selection
\fi
% Use upquote if available, for straight quotes in verbatim environments
\IfFileExists{upquote.sty}{\usepackage{upquote}}{}
\IfFileExists{microtype.sty}{% use microtype if available
  \usepackage[]{microtype}
  \UseMicrotypeSet[protrusion]{basicmath} % disable protrusion for tt fonts
}{}
\makeatletter
\@ifundefined{KOMAClassName}{% if non-KOMA class
  \IfFileExists{parskip.sty}{%
    \usepackage{parskip}
  }{% else
    \setlength{\parindent}{0pt}
    \setlength{\parskip}{6pt plus 2pt minus 1pt}}
}{% if KOMA class
  \KOMAoptions{parskip=half}}
\makeatother
\usepackage{longtable,booktabs,array}
\usepackage{calc} % for calculating minipage widths
% Correct order of tables after \paragraph or \subparagraph
\usepackage{etoolbox}
\makeatletter
\patchcmd\longtable{\par}{\if@noskipsec\mbox{}\fi\par}{}{}
\makeatother
% Allow footnotes in longtable head/foot
\IfFileExists{footnotehyper.sty}{\usepackage{footnotehyper}}{\usepackage{footnote}}
\makesavenoteenv{longtable}
\setlength{\emergencystretch}{3em} % prevent overfull lines
\providecommand{\tightlist}{%
  \setlength{\itemsep}{0pt}\setlength{\parskip}{0pt}}
\usepackage{bookmark}
\IfFileExists{xurl.sty}{\usepackage{xurl}}{} % add URL line breaks if available
\urlstyle{same}
\hypersetup{
  hidelinks,
  pdfcreator={LaTeX via pandoc}}

\author{SCX}
\date{2026}

\begin{document}

\section{Situs
Theory Final Verification Report}\label{situs-ux7406ux8bbaux6700ux7ec8ux9a8cux8bc1ux62a5ux544a}

\textbf{validate Date}: 2026-06-29\\
\textbf{validate File}: - \texttt{paper/situs\_theory/main.tex} (1518 lines) -
\texttt{theory/self\_evolution/ppe\_rigorous\_derivation.md} (1030 lines)

\begin{center}\rule{0.5\linewidth}{0.5pt}\end{center}

\subsection{Item-by-Item Check}\label{ux9010ux9879ux68c0ux67e5}

\subsubsection{1. Theorem 1.2.3 Lipschitz Constant Corrected to
2√2·π/√d？}\label{ux5b9aux7406-1.2.3-lipschitz-ux5e38ux6570ux662fux5426ux5df2ux4feeux6b63ux4e3a-22ux3c0d}

\begin{longtable}[]{@{}
  >{\raggedright\arraybackslash}p{(\linewidth - 6\tabcolsep) * \real{0.2500}}
  >{\raggedright\arraybackslash}p{(\linewidth - 6\tabcolsep) * \real{0.2500}}
  >{\raggedright\arraybackslash}p{(\linewidth - 6\tabcolsep) * \real{0.2500}}
  >{\raggedright\arraybackslash}p{(\linewidth - 6\tabcolsep) * \real{0.2500}}@{}}
\toprule\noalign{}
\begin{minipage}[b]{\linewidth}\raggedright
File
\end{minipage} & \begin{minipage}[b]{\linewidth}\raggedright
Location
\end{minipage} & \begin{minipage}[b]{\linewidth}\raggedright
Formula
\end{minipage} & \begin{minipage}[b]{\linewidth}\raggedright
Status
\end{minipage} \\
\midrule\noalign{}
\endhead
\bottomrule\noalign{}
\endlastfoot
main.tex & L405--409 &
\texttt{L\_\{\textbackslash{}PE\}\^{}\{\textbackslash{}text\{scalar\}\}\ =\ \textbackslash{}frac\{2\textbackslash{}sqrt\{2\}\textbackslash{},\textbackslash{}pi\}\{\textbackslash{}sqrt\{d\}\}\ \textbackslash{}cdot\ \textbackslash{}sqrt\{\textbackslash{}sum\_\{j=0\}\^{}\{d/2-1\}\textbackslash{}frac\{1\}\{\textbackslash{}lambda\_j\^{}2\}\}}
& ✅ \\
ppe\_rigorous\_derivation.md & L189 &
\texttt{L\_\{\textbackslash{}text\{PE\}\}\^{}\{\textbackslash{}text\{scalar\}\}\ =\ \textbackslash{}frac\{2\textbackslash{}sqrt\{2\}\textbackslash{},\textbackslash{}pi\}\{\textbackslash{}sqrt\{d\}\}\ \textbackslash{}cdot\ \textbackslash{}sqrt\{\textbackslash{}sum\_\{j=0\}\^{}\{d/2-1\}\ \textbackslash{}frac\{1\}\{\textbackslash{}lambda\_j\^{}2\}\}}
& ✅ \\
\end{longtable}

\textbf{Conclusion}: Fileone 致, number already 修正. one ization because factor  \texttt{√(2/d)}
has been incorporated into the proof derivation; the final Lipschitz constant is obtained via \texttt{√(2/d)} ×
\texttt{2π/λ\_j} sum of squares root, yielding coefficient \texttt{2√2·π/√d}. Corollaries 1.2.3
asymptotic behavior \texttt{\textasciitilde{}\ (1/ξ)·√(d/6)} thus also changes from
\texttt{O(d)} to \texttt{O(√d)}. \textbf{Pass.}

\begin{center}\rule{0.5\linewidth}{0.5pt}\end{center}

\subsubsection{2. theorem  1.2.1 one ization because factor YesNoalready 入？md 框内FormulaYesNoalready 掉
L
dependence？}\label{ux5b9aux7406-1.2.1-ux5f52ux4e00ux5316ux56e0ux5b50ux662fux5426ux5df2ux52a0ux5165md-ux6846ux5185ux516cux5f0fux662fux5426ux5df2ux53bbux6389-l-ux4f9dux8d56}

\textbf{2a. Normalization Factor √(2/d)}:

\begin{longtable}[]{@{}
  >{\raggedright\arraybackslash}p{(\linewidth - 6\tabcolsep) * \real{0.2500}}
  >{\raggedright\arraybackslash}p{(\linewidth - 6\tabcolsep) * \real{0.2500}}
  >{\raggedright\arraybackslash}p{(\linewidth - 6\tabcolsep) * \real{0.2500}}
  >{\raggedright\arraybackslash}p{(\linewidth - 6\tabcolsep) * \real{0.2500}}@{}}
\toprule\noalign{}
\begin{minipage}[b]{\linewidth}\raggedright
File
\end{minipage} & \begin{minipage}[b]{\linewidth}\raggedright
Location
\end{minipage} & \begin{minipage}[b]{\linewidth}\raggedright
Content
\end{minipage} & \begin{minipage}[b]{\linewidth}\raggedright
Status
\end{minipage} \\
\midrule\noalign{}
\endhead
\bottomrule\noalign{}
\endlastfoot
main.tex & L268--269 &
\texttt{\textbackslash{}PE\_\{\textbackslash{}text\{scalar\}\}(p,\ 2j)\ =\ \textbackslash{}sqrt\{\textbackslash{}frac\{2\}\{d\}\}\textbackslash{},\textbackslash{}sin(...)}
& ✅ \\
ppe\_rigorous\_derivation.md & L101 &
\texttt{\textbackslash{}text\{PE\}\_\{\textbackslash{}text\{scalar\}\}(p,\ 2j)\ =\ \textbackslash{}sqrt\{\textbackslash{}frac\{2\}\{d\}\}\textbackslash{},\textbackslash{}sin(...)}
& ✅ \\
\end{longtable}

Fileall 包含归one ization because factor , and all explain its using : Guarantee \texttt{‖PE(p)‖\ =\ 1}
and the encoding kernel at \texttt{d→∞} converges to the target kernel \texttt{k(Δ)} rather than diverging to
\texttt{O(d)}. 

\textbf{2b. 框内FormulaYesNo L dependence}:

\begin{longtable}[]{@{}
  >{\raggedright\arraybackslash}p{(\linewidth - 6\tabcolsep) * \real{0.2308}}
  >{\raggedright\arraybackslash}p{(\linewidth - 6\tabcolsep) * \real{0.2308}}
  >{\raggedright\arraybackslash}p{(\linewidth - 6\tabcolsep) * \real{0.2308}}
  >{\raggedright\arraybackslash}p{(\linewidth - 6\tabcolsep) * \real{0.3077}}@{}}
\toprule\noalign{}
\begin{minipage}[b]{\linewidth}\raggedright
File
\end{minipage} & \begin{minipage}[b]{\linewidth}\raggedright
Location
\end{minipage} & \begin{minipage}[b]{\linewidth}\raggedright
Formula
\end{minipage} & \begin{minipage}[b]{\linewidth}\raggedright
L dependence
\end{minipage} \\
\midrule\noalign{}
\endhead
\bottomrule\noalign{}
\endlastfoot
main.tex & L317--320 & \texttt{λ\_j\ =\ 2πξ\ ·\ cot(π(2j+1)/2d)} & No
✅ \\
ppe\_rigorous\_derivation.md & L133 &
\texttt{λ\_j\ =\ 2πξ\ ·\ cot(π(2j+1)/2d)} & No ✅ \\
\end{longtable}

L only through Corollaries 1.2.1 's Nyquist conditions \texttt{d\_min\ ≈\ L/(2ξ)}
indirect introduce , non-at code 身's框内Formulain . \textbf{Pass.}

\begin{center}\rule{0.5\linewidth}{0.5pt}\end{center}

\subsubsection{3. Theorem Numbering Unified to
4.1？}\label{ux5b9aux7406ux7f16ux53f7ux662fux5426ux7edfux4e00ux4e3a-4.1}

``theorem  4.1'' 指'sYes''Learningtype  PE break  Theorem 3 'sprecise conditions''this theorem . 

\begin{longtable}[]{@{}llll@{}}
\toprule\noalign{}
File & Location & 编号 & target  \\
\midrule\noalign{}
\endhead
\bottomrule\noalign{}
\endlastfoot
main.tex & L1095 & theorem  4.1 &
\texttt{\textbackslash{}label\{thm:4.1\}} \\
ppe\_rigorous\_derivation.md & L775 & theorem  4.1 & --- \\
\end{longtable}

Fileone make using  \textbf{theorem  4.1}. Meanwhile Proposition 4.1(fixed  PE
maintain non-can be property )numbering is also consistent.\textbf{Pass.}

\begin{center}\rule{0.5\linewidth}{0.5pt}\end{center}

\subsubsection{4. d\_min Corrected to
L/(2ξ)？}\label{d_min-ux662fux5426ux4feeux6b63ux4e3a-l2ux3be}

\begin{longtable}[]{@{}
  >{\raggedright\arraybackslash}p{(\linewidth - 6\tabcolsep) * \real{0.2500}}
  >{\raggedright\arraybackslash}p{(\linewidth - 6\tabcolsep) * \real{0.2500}}
  >{\raggedright\arraybackslash}p{(\linewidth - 6\tabcolsep) * \real{0.2500}}
  >{\raggedright\arraybackslash}p{(\linewidth - 6\tabcolsep) * \real{0.2500}}@{}}
\toprule\noalign{}
\begin{minipage}[b]{\linewidth}\raggedright
File
\end{minipage} & \begin{minipage}[b]{\linewidth}\raggedright
Location
\end{minipage} & \begin{minipage}[b]{\linewidth}\raggedright
Formula
\end{minipage} & \begin{minipage}[b]{\linewidth}\raggedright
推导
\end{minipage} \\
\midrule\noalign{}
\endhead
\bottomrule\noalign{}
\endlastfoot
main.tex & L368 &
\texttt{d\_\{\textbackslash{}min\}\ \textbackslash{}approx\ \textbackslash{}frac\{L\}\{2\textbackslash{}xi\}}
& λ\_max ≈ 4dξ, Nyquist: λ\_max ≥ 2L → 4dξ ≥ 2L → d ≥ L/(2ξ) \\
ppe\_rigorous\_derivation.md & L173 &
\texttt{d\_\{\textbackslash{}min\}\ \textbackslash{}approx\ \textbackslash{}frac\{L\}\{2\textbackslash{}xi\}}
& on  \\
\end{longtable}

推导链Complete, Nyquist conditionsalready Correctusing (覆盖whole order sequence 区间requires  λ\_max ≥
2L). \textbf{Pass.}

\begin{center}\rule{0.5\linewidth}{0.5pt}\end{center}

\subsubsection{5. Theorem 2.4.1
Downgraded to Heuristic Estimate?}\label{ux5b9aux7406-2.4.1-ux662fux5426ux5df2ux964dux7ea7ux4e3aux542fux53d1ux5f0fux4f30ux8ba1}

\textbf{Key Change}:
From''Theorem (lower bound)``→''Proposition (Heuristic Estimate---Not a Rigorous lower bound)``. 

\begin{longtable}[]{@{}
  >{\raggedright\arraybackslash}p{(\linewidth - 4\tabcolsep) * \real{0.1395}}
  >{\raggedright\arraybackslash}p{(\linewidth - 4\tabcolsep) * \real{0.2326}}
  >{\raggedright\arraybackslash}p{(\linewidth - 4\tabcolsep) * \real{0.6279}}@{}}
\toprule\noalign{}
\begin{minipage}[b]{\linewidth}\raggedright
feature
\end{minipage} & \begin{minipage}[b]{\linewidth}\raggedright
main.tex
\end{minipage} & \begin{minipage}[b]{\linewidth}\raggedright
ppe\_rigorous\_derivation.md
\end{minipage} \\
\midrule\noalign{}
\endhead
\bottomrule\noalign{}
\endlastfoot
环境class type  & \texttt{\textbackslash{}begin\{proposition\}} &
\texttt{**命题\ 2.4.1**} \\
诚target  & \texttt{\textbackslash{}heuristic\{\}} + 警告框 &
\texttt{{[}formula {]}} + 警告框 \\
明确声明 & ``此表达formula \textbf{non-Yes}严格'snumber learning lower bound'' & ``此表达formula 
\emph{non-Yes} 严格'snumber learning lower bound'' \\
逻辑Errortarget  & ``FromFanolower boundnon-can 逻辑推差'slower bound'' & ``From Fano
lower bound\ldots non-can 逻辑推'' \\
etc.号误using Explanation &
``if Fanonon-etc.formula at bound all to etc.号(this at one 般况under non-立)'' & ``if Assumptions
Fano non-etc.formula at bound all to etc.号(requires Distributionconditions)'' \\
\end{longtable}

Filefor 齐良. \textbf{Pass.}

\begin{center}\rule{0.5\linewidth}{0.5pt}\end{center}

\subsubsection{6. Theorem 2.2.1 Step 4
Restricted to Bayes Optimal?}\label{ux5b9aux7406-2.2.1-ux6b65ux9aa4-4-ux662fux5426ux5df2ux9650ux5b9aux4e3aux8d1dux53f6ux65afux6700ux4f18}

\begin{longtable}[]{@{}
  >{\raggedright\arraybackslash}p{(\linewidth - 4\tabcolsep) * \real{0.1395}}
  >{\raggedright\arraybackslash}p{(\linewidth - 4\tabcolsep) * \real{0.2326}}
  >{\raggedright\arraybackslash}p{(\linewidth - 4\tabcolsep) * \real{0.6279}}@{}}
\toprule\noalign{}
\begin{minipage}[b]{\linewidth}\raggedright
level 
\end{minipage} & \begin{minipage}[b]{\linewidth}\raggedright
main.tex
\end{minipage} & \begin{minipage}[b]{\linewidth}\raggedright
ppe\_rigorous\_derivation.md
\end{minipage} \\
\midrule\noalign{}
\endhead
\bottomrule\noalign{}
\endlastfoot
theorem target  & ``充conditions------Information Theoryformula (贝叶斯Optimal限)'' &
``充conditions------Information Theoryformula , 贝叶斯Optimal限'' \\
Conclusion陈述 & ``\textbf{贝叶斯Optimalclass device }at 增广featureunder 'sDetection边际变ization satisfies 
δ\textgreater0'' &
``\textbf{贝叶斯Optimalclass device }at 增广featureunder 'sDetection边际变ization satisfies 
δ\textgreater0'' \\
步骤 4 target  & \texttt{\textbackslash{}heuristic\{\}} + ``猜想'' &
\texttt{{[}formula {]}} + ``猜想'' \\
局限声明 & ``步骤1-3as 严格Information Theory推导(Conclusionfor 贝叶斯Optimalclass device 立)'' &
``步骤 1-3 Proof贝叶斯Optimalclass device can be with FromLocation信息in 受益'' \\
退ization 声明 & ``if 步骤4'sgeneralize non-立, theorem ization as via Conclusion------平凡's'' &
``if 步骤 4 'sgeneralize non-立, theorem  2.2.1 ization as one via Conclusion------平凡's'' \\
\end{longtable}

\textbf{Pass.}

\begin{center}\rule{0.5\linewidth}{0.5pt}\end{center}

\subsubsection{7. Theorem 3'
break toward YesNoalready as Problem？}\label{theorem-3-ux7834ux574fux65b9ux5411ux662fux5426ux5df2ux6539ux4e3aux5f00ux653eux95eeux9898}

\begin{longtable}[]{@{}
  >{\raggedright\arraybackslash}p{(\linewidth - 4\tabcolsep) * \real{0.1395}}
  >{\raggedright\arraybackslash}p{(\linewidth - 4\tabcolsep) * \real{0.2326}}
  >{\raggedright\arraybackslash}p{(\linewidth - 4\tabcolsep) * \real{0.6279}}@{}}
\toprule\noalign{}
\begin{minipage}[b]{\linewidth}\raggedright
Location
\end{minipage} & \begin{minipage}[b]{\linewidth}\raggedright
main.tex
\end{minipage} & \begin{minipage}[b]{\linewidth}\raggedright
ppe\_rigorous\_derivation.md
\end{minipage} \\
\midrule\noalign{}
\endhead
\bottomrule\noalign{}
\endlastfoot
Theorem 3' 陈述内 & ``if 前提 (A) non-立:
non-can be property YesNois break Purpose: is a\textbf{放Problem}'' & ``if 前提 (A) non-立:
non-can be property YesNois break Purpose: is a\textbf{放Problem}'' \\
诚target  & \texttt{\textbackslash{}openproblem\{\}} &
\texttt{{[}Problem{]}} \\
minimum e.g.重评 & ``此示e.g.际on ProofTheorem 3's鲁棒property , while non-its 脆弱property '' &
``此minimum e.g.际on Proof Theorem 3 's鲁棒property , while non-its 脆弱property '' \\
诚表 & theorem 4.1Purpose: :
\texttt{\textbackslash{}openproblem\ break toward Noconstruct property Proof} & --- \\
Theorem 3'条Purpose:  &
\texttt{\textbackslash{}rigorous\ 前提立;\ \textbackslash{}openproblem\ 前提non-立}
& --- \\
\end{longtable}

原声称's''construct property e.g.Proofbreak toward ''already is 重评as ''construct property e.g.反while theorem 's鲁棒property ''------符预期. \textbf{Pass.}

\begin{center}\rule{0.5\linewidth}{0.5pt}\end{center}

\subsubsection{8.
份FileYesNoalso has Symbolnon-one 致？}\label{ux4e24ux4efdux6587ux4ef6ux4e4bux95f4ux662fux5426ux8fd8ux6709ux7b26ux53f7ux4e0dux4e00ux81f4}

\paragraph{8a. Consistent Items ✅}\label{a.-ux4e00ux81f4ux7684ux9879ux76ee}

\begin{longtable}[]{@{}
  >{\raggedright\arraybackslash}p{(\linewidth - 6\tabcolsep) * \real{0.2037}}
  >{\raggedright\arraybackslash}p{(\linewidth - 6\tabcolsep) * \real{0.1852}}
  >{\raggedright\arraybackslash}p{(\linewidth - 6\tabcolsep) * \real{0.5000}}
  >{\raggedright\arraybackslash}p{(\linewidth - 6\tabcolsep) * \real{0.1111}}@{}}
\toprule\noalign{}
\begin{minipage}[b]{\linewidth}\raggedright
Symbol/概念
\end{minipage} & \begin{minipage}[b]{\linewidth}\raggedright
main.tex
\end{minipage} & \begin{minipage}[b]{\linewidth}\raggedright
ppe\_rigorous\_derivation.md
\end{minipage} & \begin{minipage}[b]{\linewidth}\raggedright
Status
\end{minipage} \\
\midrule\noalign{}
\endhead
\bottomrule\noalign{}
\endlastfoot
Detection边际 & \texttt{Δ\_s\ =\ p\_noisy\ -\ p\_clean} &
\texttt{Δ\_s\ =\ p\_noisy\ -\ p\_clean} & ✅ \\
PPE增强边际 & \texttt{Δ\_s\^{}\{Situs\}\ =\ Δ\_s\ +\ δ\_s\^{}\{PE\}} &
\texttt{Δ\_s\^{}\{PPE\}\ =\ Δ\_s\ +\ δ\_s\^{}\{PE\}} & ✅ \\
Chernoff-Hoeffding & \texttt{exp(-2M(Δ\_s\ +\ δ\_s\^{}\{PE\})\^{}2)} &
\texttt{exp(-2M(Δ\_s\ +\ δ\_s\^{}\{PE\})\^{}2)} & ✅ \\
编code non-完美 &
\texttt{ε\_PE\ =\ I(Y;P\textbackslash{}\textbar{}X)\ -\ I(Y;PE(P)\textbackslash{}\textbar{}X)}
&
\texttt{ε\_PE\ =\ I(Y;P\textbackslash{}\textbar{}X)\ -\ I(Y;PE(P)\textbackslash{}\textbar{}X)}
& ✅ \\
Theorem 2' Upper Bound & \texttt{F1\_base\ +\ C\_F·√((δ\ +\ 2ε/C\_F²)/2)} &
\texttt{F1\_base\ +\ C\_F·√((δ\ +\ 2ε/C\_F²)/2)} & ✅ \\
d\_min & \texttt{L/(2ξ)} & \texttt{L/(2ξ)} & ✅ \\
target quantity PEdefinition  & \texttt{√(2/d)·sin/cos(2πp/λ\_j)} &
\texttt{√(2/d)·sin/cos(2πp/λ\_j)} & ✅ \\
3D旋转PE & \texttt{R(p)·e\_0} & \texttt{R(p)·e\_0} & ✅ \\
旋转Lipschitz & \texttt{max(α,\ β,\ γ)} & \texttt{max(α,\ β,\ γ)} &
✅ \\
正弦Lipschitz & \texttt{(2√2·π/√d)·√(Σ\ 1/λ\_j²)} &
\texttt{(2√2·π/√d)·√(Σ\ 1/λ\_j²)} & ✅ \\
theorem 4.1conditions & \texttt{L\_total\ =\ L\_sup\ +\ λ·L\_aux} &
\texttt{mathcal\{L\}\_\{total\}\ =\ mathcal\{L\}\_\{sup\}\ +\ λ·mathcal\{L\}\_\{aux\}}
& ✅ \\
\end{longtable}

\paragraph{8b.
预期Difference(non-Problem)}\label{b.-ux9884ux671fux5deeux5f02ux975eux95eeux9898}

\begin{longtable}[]{@{}
  >{\raggedright\arraybackslash}p{(\linewidth - 4\tabcolsep) * \real{0.3333}}
  >{\raggedright\arraybackslash}p{(\linewidth - 4\tabcolsep) * \real{0.3333}}
  >{\raggedright\arraybackslash}p{(\linewidth - 4\tabcolsep) * \real{0.3333}}@{}}
\toprule\noalign{}
\begin{minipage}[b]{\linewidth}\raggedright
Difference
\end{minipage} & \begin{minipage}[b]{\linewidth}\raggedright
Reason
\end{minipage} & \begin{minipage}[b]{\linewidth}\raggedright
Judgment
\end{minipage} \\
\midrule\noalign{}
\endhead
\bottomrule\noalign{}
\endlastfoot
\texttt{\textbackslash{}Situs\{\}} vs \texttt{PPE} & paper make using framework 名
Situs, derivation make using 缩写 PPE & 预期 ✅ \\
\texttt{\^{}\{\textbackslash{}Situs\}} vs
\texttt{\^{}\{\textbackslash{}text\{PPE\}\}} & on target non- & 预期 ✅ \\
章节编号偏移 & paper has  §1 引言, derivation No引言direct Fromcode definition  &
预期 ✅ \\
Theorem 环境形formula  & LaTeX \texttt{\textbackslash{}begin\{theorem\}} vs
Markdown \texttt{**theorem **} & formula Difference ✅ \\
\texttt{\textbackslash{}PE} vs \texttt{PE} & LaTeX operator vs plain
text & formula Difference ✅ \\
\end{longtable}

\paragraph{8c. ⚠️
Residual Inconsistencies (Require Attention)}\label{c.-ux6b8bux4f59ux4e0dux4e00ux81f4ux9700ux5173ux6ce8}

\textbf{R1. Abstract in 's''充must conditions''声称}

\begin{longtable}[]{@{}
  >{\raggedright\arraybackslash}p{(\linewidth - 2\tabcolsep) * \real{0.5000}}
  >{\raggedright\arraybackslash}p{(\linewidth - 2\tabcolsep) * \real{0.5000}}@{}}
\toprule\noalign{}
\begin{minipage}[b]{\linewidth}\raggedright
Location
\end{minipage} & \begin{minipage}[b]{\linewidth}\raggedright
Content
\end{minipage} \\
\midrule\noalign{}
\endhead
\bottomrule\noalign{}
\endlastfoot
main.tex L97 (摘must ) & ``Proof \(\delta_s^{\PE} > 0\)
's\textbf{充must conditions}as  \(I(Y; P \mid S) > 0\)(theorem 2.2.1)'' \\
main.tex L649 (theorem  2.2.1 target 题) & ``\(\delta_s^{\PE} > 0\)
's\textbf{充conditions}------Information Theoryformula (贝叶斯Optimal限)'' \\
main.tex L694 (命题 2.2.1 target 题) & ``\(\delta_s^{\PE} > 0\)
's\textbf{必must conditions}'' \\
ppe\_rigorous\_derivation.md L1020 (Summary) & ``Proof
\(\delta_s^{\PE} > 0\)
's\textbf{充conditions}(\(I(Y; P \mid S) > 0\))'' \\
\end{longtable}

\textbf{Problem}: must theorem  2.2.1 give ''充must conditions'', but theorem  2.2.1
only 证''充conditions''(and at 贝叶斯Optimalclass device ), must conditionsby \textbf{independent 's}命题
2.2.1 additionally give (and 允许non-Information TheoryMechanism). 此, conditions包含步骤 4
'sformula generalize (From贝叶斯Optimalto expert ), while non-纯粹严格Proof. 

\textbf{建议}:
摘must in ''充must conditions''should as ''充conditions(theorem 2.2.1)and must conditions(命题2.2.1)``, or to will 归属From单独's''theorem 2.2.1''修正as ''theorem 2.2.1
+ 命题2.2.1''. 

\textbf{R2. ppe\_rigorous\_derivation.md theorem  2.3.1 's双重陈述}

\begin{longtable}[]{@{}
  >{\raggedright\arraybackslash}p{(\linewidth - 2\tabcolsep) * \real{0.5000}}
  >{\raggedright\arraybackslash}p{(\linewidth - 2\tabcolsep) * \real{0.5000}}@{}}
\toprule\noalign{}
\begin{minipage}[b]{\linewidth}\raggedright
Location
\end{minipage} & \begin{minipage}[b]{\linewidth}\raggedright
Content
\end{minipage} \\
\midrule\noalign{}
\endhead
\bottomrule\noalign{}
\endlastfoot
L458 &
粗略形formula : \texttt{min(1,\ √(½·I(P;PE(P)\textbackslash{}\textbar{}S,Y)))\ +\ Bernstein修正项} \\
L462 &
precise formula : \texttt{min(1-p\_clean,\ √(½·D\_KL(P\_\{PE\textbackslash{}\textbar{}S,Y\}\ ‖\ P\_\{PE\textbackslash{}\textbar{}S\})))\ +\ δ\_s\^{}\{variance\}} \\
main.tex L723 & only has precise formula  ✅ \\
\end{longtable}

\textbf{Problem}: ppe\_rigorous\_derivation.md
先give one 较粗略'sformula , again give precise formula . 版'sPartnon-(\texttt{I(P;PE\textbar{}S,Y)}
vs \texttt{D\_KL}), and Partone make using 概念on non-'s
\texttt{Bernstein\ 修正项} while precise make using 
\texttt{δ\_s\^{}\{variance\}}. this can be can 导致读者困惑------形formula non-etc.价. 

\textbf{建议}: ppe\_rigorous\_derivation.md
should at 粗略形formula 前''粗略''target 记, Explanationprecise formula under . or direct 删除粗略形formula , only 保留and 
main.tex one 致'sprecise 版. 

\begin{center}\rule{0.5\linewidth}{0.5pt}\end{center}

\subsection{Overall Assessment}\label{ux603bux4f53ux8bc4ux4f30}

\begin{longtable}[]{@{}ll@{}}
\toprule\noalign{}
Check Item & Status \\
\midrule\noalign{}
\endhead
\bottomrule\noalign{}
\endlastfoot
1. Lipschitz number 修正 & ✅ through  \\
2. one ization because factor  +  L dependence & ✅ through  \\
3. theorem 编号统one as  4.1 & ✅ through  \\
4. d\_min = L/(2ξ) & ✅ through  \\
5. Theorem 2.4.1 level as formula  & ✅ through  \\
6. Theorem 2.2.1 Step 4 贝叶斯Optimal & ✅ through  \\
7. Theorem 3' break toward  → Problem & ✅ through  \\
8. FileSymbolone property  & ⚠️ 2 项残余Problem(见 R1, R2) \\
\end{longtable}

\textbf{Summary}: 7/8 complete Pass.core 修正(Lipschitz
number , one ization because factor , Fano level , d\_min
修正, 贝叶斯Optimal限, break toward Problemization )at Filein all already Correctand maintain one 致. 2
项残余Problemas  minor level 别------Abstract wording too absolute, and derivation
dual theorem statement in the document------do not affect mathematical correctness and can be corrected in the next editing round.

\end{document}
